\documentclass[openany,fontset=fandol]{ctexbook}
\usepackage{amsmath,amssymb}
\usepackage[unicode,hidelinks]{hyperref}
\newcommand{\note}[1]{\footnote{译注：#1}}
\begin{document}
\title{代数几何基础第五卷预备笔记\\定理13.2：初等邻域的构造}
\author{据用户上传英译编辑稿译出}
\date{2026年11月17日}
\maketitle
\chapter*{范围与校订状态}
本批译出跨底本第50—51页的一长段，从“There exists a linar subvariety $C$ of $L$”至“This is the desired neighborhood of $x$ contained in $X-Y$”。为便于阅读，按构造步骤分段。下一接点为第51页“It remains only to prove Lemma 13.3”。

原段多处对象和字母错置可由同段前后式辨明，正文整理并在脚注存原写法；关于横截选择、经过$x$与连通性的论证略写明确保留，未把引理13.3的证明标为完成。本稿未覆盖主稿。
\chapter*{§13 初等态射与M. Artin定理（续）}
存在$L_0$中的一个线性子簇$C$，它在$\mathbf P^r$中的余维为$n$，即在$L_0$中的余维为1，并且不与有限集$(L_0\cap Y)\cup\{x\}$相交。令$T=X\cap C$，则$T$是$X$中一个在$k$上平展的非空子概形，不含$x$，且与$Y$不交。\note{底本两次写$L$，同段的$L_0\cap Y$及随后参数点所代表的子空间明确指向引理中的$L_0$，故统一为$L_0$。但仅仅避开所列有限集，并不足以从字面上推出$X\cap C$平展；还需适当横截选择$C$。此处保留原作者对所选$C$的性质的断言，并明确指出该选择步骤未在本段展开，不把任意满足避开条件的$C$都说成适用。}

另一方面，考虑$\operatorname{Grass}^{n-1}(P)$中对应于包含$C$的线性子簇的子概形$Q$。$Q$是一个维数为$n-1$的射影空间；特别地，它在$k$上光滑，维数为$n-1$。$L_0$对应于一个点$\xi_0\in Q(k)$。沿用此前引入的一般记号，考虑$X^{(n-1)}$沿浸入$Q\to\operatorname{Grass}^{n-1}(P)$的逆像$X_Q^{(n-1)}$，以及$Y_Q^{(n-1)}$、$T_Q^{(n-1)}$。后二者也是$X_Q^{(n-1)}$中互不相交的闭子概形。以$p,q,r$分别记这三个概形到$Q$的结构投影。\note{底本先写参数点$\xi$，下页续文改写$\xi_0$，此处统一为$\xi_0$。列举第三个逆像时底本只印出带上下标的$Q$而漏掉对象字母；紧随其后的$T_Q^{n-1}\simeq T\times_k Q$及边界构造确定它是$T$的逆像。括号中的上标$(n-1)$表示余维，排版整理不改变原含义。}

由假设，$p$在$\xi_0$上方的点处光滑，$q$在$\xi_0$上方的点处平展；$r$也是如此，因为$T_Q^{(n-1)}$无非就是$T\times_k Q$（在$Q$上同构）。最后，态射$p$固有；又由于$X$几何连通，$p$的纤维几何连通（Bertini定理）。\note{这一连通性结论按原稿引用保留。底本并未在此写明所用Bertini定理的完整假设或引用位置，本批不以“$X$几何连通”单独替代所需的线性截面连通性论证。}

\newpage
因此，存在$\xi_0$在$Q$中的一个开邻域$V$，使得
\[
X'=X_Q^{(n-1)}\big|_V
\]
在$V$上固有、光滑，并有几何连通纤维。由于$\xi_0$处的纤维正是$X\cap L_0$，其维数为1，可以假设$X'$在$V$上的所有纤维都为一维。最后，将$V$取得充分小，可以假设$Y_Q^{(n-1)}|_V$及$T_Q^{(n-1)}|_V$在$V$上平展，因而二者的和概形$Z'$（它可识别为$X'$的一个闭子概形）在$V$上平展。于是，令
\[
U=X'\setminus Z',
\]
便得到一个初等态射$U\to V$。\note{底本原句为“an open neighborhood $U$ of $Q$ in $X$”，但紧接着把族限制到$V$，上下文始终是在参数点附近缩小基底，故正文订为“$\xi_0$在$Q$中的开邻域$V$”；原“fiber of $Q$”同理订为$\xi_0$处的纤维。定义13.1所需射影性可由原射影关联族限制得到；边界的满射性和恒定次数还须使用非空$T$及必要的缩小。本段原文没有把这些核对全部写出。}

但$U$也在
\[
X''=X_Q^{(n-1)}\setminus
\bigl(Y_Q^{(n-1)}\cup T_Q^{(n-1)}\bigr)
\]
中开；$X''$是$X\setminus(Y\cup T)$在$X_Q^{(n-1)}$中的逆像。另一方面，$X''\to X\setminus(Y\cup T)$是同构，因为
\[
X_Q^{(n-1)}\setminus T_Q^{(n-1)}\longrightarrow X\setminus T
\]
是同构。因此，$U$可识别为$X\setminus(Y\cup T)$的一个开子集；它包含$(L_0\cap X)\setminus(Y\cup T)$，因而包含$x$。这就是所需的、包含在$X\setminus Y$中的$x$的邻域。\note{底本在$X''$的公式中写$Z_Q^{n-1}$，但随即说明它是$X-Y-T$的逆像，且本段移去的第二个边界一直是$T$，故订为$T_Q^{(n-1)}$。底本又写$U$包含整个$L_0\cap X$；因$U$已删去边界，这应为删去$Y\cup T$后的部分，正文据构造修正。最后“因而包含$x$”还使用$x\in L_0$：底本引理13.3的陈述漏明这一点，下一段证明却明确取经过$x$的线性子簇。此项前后依赖在这里指出，不能只凭上一批引理字面推出。}
\end{document}
